\textbf{Proof.}
The entries of \(S_B=B^{\top}\Omega B\) are quadratic polynomials in the \(2n\) coordinates of \((b_1,b_0)\), with rational or real-algebraic coefficients.  The two sphere conditions are polynomial equalities, while rational box and signed pivot-chart restrictions are polynomial weak inequalities.  For a symmetric two-by-two matrix,
\[
 \begin{bmatrix}u&v\\v&z\end{bmatrix}\succeq0
 \quad\Longleftrightarrow\quad
 u\ge0,\ z\ge0,\ uz-v^2\ge0.
\]
Applying this identity to \(S_B-n\kappa I_2\) includes equality at the regularity boundary exactly.

Thus chart-box feasibility is an existential first-order formula over the real closed field of real-algebraic numbers.  The Collins EVAL/ELIM result supplies a quantifier-free equivalent.  Its DECOMP construction partitions the finitely many polynomial sign conditions into finitely described cells with algebraic sample points; recursive projection lowers the variable count and therefore terminates after finitely many stages.  Exact algebraic-number isolation evaluates every sign at each sample.  A sign-compatible sample satisfying the formula is an exact feasible point; if no cell sample satisfies it, the finite sign-invariant decomposition is an infeasibility certificate.  These are the adaptations to the present sphere, box, and positive-semidefinite constraints; none is asserted by the cited leaf alone. \(\square\)